\documentclass[10pt,a4paper]{amsart}
\usepackage[margin=19mm]{geometry}
\usepackage{amsmath,amssymb,mathtools}
\usepackage[scr=rsfs,cal=euler]{mathalfa}
\usepackage{xfrac}
\usepackage[unicode,hidelinks,bookmarks=false]{hyperref}
\newcommand{\ie}{i.e.\ }
\newcommand{\cf}{cf.\ }
\newcommand{\resp}{respectively\ }
\newcommand{\Subsection}{\subsection}
\providecommand{\nobreakdash}{\nobreak\hbox{-}}
\setlength{\emergencystretch}{2em}
\multlinegap0pt
\providecommand{\skpt}{\par\medskip}
\newcommand\mto{\mathchoice{\longmapsto}{\mapsto}{\mapsto}{\mapsto}}
\newcommand\hto{\mathrel{\lhook\joinrel\to}}
\newcommand{\onto}{\mathchoice{\to\hspace*{-5mm}\to}{\to\hspace*{-2.85mm}\to}{\to\hspace*{-3mm}\to}{\to\hspace*{-3mm}\to}}

\newcommand{\sbullet}{{\scriptscriptstyle\bullet}}
\newcommand{\csbullet}{{\raisebox{1pt}{$\sbullet$}}}

\let\ts\textstyle

\usepackage{tikz}
\usetikzlibrary{cd}
\tikzcdset{arrow style=Latin Modern}
\newcommand{\dradpl}{{\hspace*{-2.5mm}\begin{tikzcd}[ampersand replacement=\&,column sep = scriptsize]\mbox{}\ar[r,dashed]\&\mbox{}\end{tikzcd}\hspace*{-2.5mm}}}
\newcommand{\dratst}{{\hspace*{-2.5mm}\begin{tikzcd}[ampersand replacement=\&,column sep = small]\mbox{}\ar[r,dashed]\&\mbox{}\end{tikzcd}\hspace*{-2mm}}}
\let\drascst\dratst
\newcommand{\dra}{\mathchoice{\dradpl}{\dratst}{\drascst}{\drascst}}
\let \dashrightarrow\dra

\let\cal\mathcal
\let\bb\mathbb\let\Bbb\mathbb
\let\mf\mathfrak
\usepackage{tikz-cd}
\usepackage{mathtools}

\theoremstyle{plain}
\newtheorem{thm}{Theorem}[section]
\newtheorem{lem}[thm]{Lemma}
\newtheorem{prop}[thm]{Proposition}
\newtheorem{cor}[thm]{Corollary}

\theoremstyle{definition}
\newtheorem{defn}[thm]{Definition}
\newtheorem{exmp}[thm]{Example}
\newtheorem{rem}[thm]{Remark}
\newtheorem{rems}[thm]{Remarks}

\DeclareMathOperator{\bAlb}{\mathbf{Alb}}
\DeclareMathOperator{\alg}{alg}
\DeclareMathOperator{\bl}{bl}
\DeclareMathOperator{\Br}{Br}
\DeclareMathOperator{\CH}{CH}
\DeclareMathOperator{\bCH}{\mathbf{CH}}
\DeclareMathOperator{\Gal}{Gal}
\DeclareMathOperator{\fppf}{fppf}
\DeclareMathOperator{\Hilb}{Hilb}
\DeclareMathOperator{\id}{id}
\DeclareMathOperator\Ker{Ker}
\DeclareMathOperator{\NS}{NS}
\DeclareMathOperator{\Pic}{Pic}
\DeclareMathOperator{\bPic}{\mathbf{Pic}}
\DeclareMathOperator{\Proj}{Proj}
\DeclareMathOperator{\Sym}{Sym}

\newcommand{\defi}[1]{\emph{#1}}
\newcommand{\kbar}{{\overline{k}}}

\let\phi\varphi

\renewcommand\P{\mathbb{P}}
\renewcommand\O{\mathcal{O}}
\newcommand\Z{\mathbb{Z}}
\newcommand\C{\mathbb{C}}
\newcommand\Q{\mathbb{Q}}
\newcommand\R{\mathbb{R}}
\newcommand\F{\mathbb{F}}

\title{Arithmetic and birational properties of linear spaces on intersections of two quadrics: Theorem 1.3}
\author{Lena Ji and Fumiaki Suzuki}
\date{Source: J.\'E.P.--M. 12 (2025), 1161--1196; DOI 10.5802/jep.308}
\begin{document}\maketitle
\noindent\textit{Editorial scope.} The counted unit is the original Theorem 1.3, with its complete original statement and proof including the stated zero-dimensional alternative. The original local quadratic-form, hyperbolic-reduction and incidence constructions used by that proof are retained as uncounted core. Later arithmetic applications, Chow-scheme theory and real-isotopy classification are omitted. Original theorem and lemma numbering and actual published bibliography labels are retained.
\section{Original selected result and definitions}
To state this result, we~first need to introduce a definition. If~$F_{r}(X)(k)\neq\emptyset$, choose $\ell \in F_r(X)(k)$ and define $\mathcal{Q}^{(r)}\to \P^1$ to be the hyperbolic reduction of $\mathcal{Q}\to \P^1$ with respect to $\ell$ (see Section~\ref{subsection:hyperbolic-reduction}).
The hyperbolic reduction $\mathcal{Q}^{(r)}\to \P^1$ is itself a quadric fibration
and may be regarded as the relative Fano scheme of isotropic $(r+1)$-planes of $\mathcal{Q}\to \P^1$ containing $\ell$;
moreover, the $k$\nobreakdash-birational equivalence class of $\mathcal{Q}^{(r)}$ does not depend on $\ell$ (see Section~\ref{sec:Q^(r)-construction} for these and additional properties). We~prove:

\setcounter{thm}{2}
\begin{thm}\label{thm:symmetric}
Over a field $k$ of characteristic $\neq 2$, let $X$ be a smooth complete intersection of two quadrics in $\P^N$.
Let $r$ be such that $0\leq r\leq \lfloor \sfrac{N}{2}\rfloor -1$. If~$F_r(X)(k)\neq \emptyset$, then one of the following conditions holds.
\begin{enumerate}
\item\label{thm:symmetric1} $F_r(X)$ is $k$\nobreakdash-birational to $\Sym^{r+1}\mathcal{Q}^{(r)}$.
\item\label{thm:symmetric2} $N=2g$ and $r=g-1$ for some $g\geq 1$. In~this case,
the subscheme of $F_{g-1}(X)$ parametrizing $(g-1)$-planes on $X$
disjoint from $\ell$ is $k$\nobreakdash-isomorphic to
the subscheme of $\Sym^{g}\mathcal{Q}^{(g-1)}$ parametrizing $g$-tuples of distinct points of $\mathcal{Q}^{(g-1)}$, and they are \(0\)\nobreakdash-dimensional schemes of length $\binom{2g+1}{g}$.
\end{enumerate}
\end{thm}

\section{Preliminaries}\label{sec:background}

\Subsection{Fano schemes of $r$-planes on a smooth complete intersection of two quadrics}

Let \(X\subset\bb P^N\) be a smooth complete intersection of two quadrics. For non-negative integers $r\geq 0$, $F_r(X)$ denotes the Fano scheme of $r$-planes on $X$. In~this section, we~recall some preliminary results about these Fano schemes \(F_r(X)\).

\begin{lem}[{\cite[Lem.\,A.3 \& A.4]{CT-intersection-quadrics}}]\label{lem:Kuznetsov_fanoscheme}
Over a field \(k\) of characteristic \(\neq 2\), let $X$ be a smooth complete intersection of two quadrics in $\P^{N}$. The following hold.
\begin{enumerate}
\item
If $r>\lfloor\sfrac{N}{2}\rfloor -1$,
$F_r(X)$ is empty.
\item
If $0\leq r\leq \lfloor\sfrac{N}{2}\rfloor -1$,
$F_r(X)$ is non-empty, smooth, projective, and of dimension $(r+1)(N-2r-2)$.
\item If $0\leq r< \sfrac{N}{2} -1$ (or equivalently if $\dim F_r(X)>0$), then $F_r(X)$ is geometrically connected.
\end{enumerate}
\end{lem}

\setcounter{thm}{2}
\subsection{Quadric bundles and hyperbolic reductions}\label{subsection:hyperbolic-reduction}

In this section, we~recall the definition and basic properties of quadric fibrations and of the hyperbolic reductions of a quadric fibrations; see, e.g., \cite[\S 2]{Kuznetsov-hyperbolic}, \cite[\S 1.3]{ABB14} for more details.

Let \(S\) be an integral separated Noetherian scheme over a field \(k\) of characteristic \(\neq 2\). A \defi{quadric fibration} over \(S\) is a morphism \(\cal Q\to S\) that can be written as a composition \(\cal Q\hto\bb P_S(\cal E) \coloneqq \underline{\Proj}_S \Sym^\csbullet (\cal E^\vee) \to S\)
where \(\cal E\) is a vector bundle and \(\cal Q\hto\bb P_S(\cal E)\) is a divisor of relative degree 2 over \(S\). A quadric fibration is determined by a quadratic form \(q\colon\Sym^2\cal E\to\cal L^\vee\) with values in a line bundle \(\cal L\), and we define the \defi{degeneration divisor} of the quadric fibration to be the zero locus of the determinant of $q$.

\begin{lem}[{\cite[Prop.\,1.2.5]{ABB14}}]\label{lem:quadric-bundle-simple-degeneration}
Let \(S\) be a smooth scheme over a field of characteristic \(\neq 2\), and let \(\pi\colon\cal Q\to S\) be a flat quadric fibration with smooth generic fiber. Then the degeneration divisor of \(\pi\) is smooth over \(k\) if and only if \(\cal Q\) is smooth over \(k\) and \(\pi\) has simple degeneration (i.e., the fibers of \(\pi\) have corank at most $1$).
\end{lem}

A subbundle \(\cal F\subset \cal E\) is \defi{isotropic} if \(q|_{\cal F}=0\) (or, equivalently, if \(\bb P_S(\cal F)\subset\cal Q\)). An~isotropic subbundle \(\cal F\subset\cal E\) is \defi{regular} if, for every (closed) point \(s\in S\), the fiber \(\bb P_S(\cal F)_s\) is contained in the smooth locus of the fiber \(\cal Q_s\). If~\(\cal F\) is regular isotropic, then~\(\cal F\) is contained in the subbundle\vspace*{-3pt}\enlargethispage{.5\baselineskip}
\[
\cal F^\perp\coloneqq\Ker(\cal E\onto\cal F^\vee\otimes\cal L^\vee)
\]
of \(\cal E\). Moreover, \(\cal F\) is in the kernel of the restriction \(q|_{\cal F^\perp}\), so we have an induced quadratic form on \(\cal F^\perp/\cal F\).

\begin{defn}
The induced quadratic form \(\overline{q}\colon\Sym^2(\cal F^\perp/\cal F)\to\cal L^{\vee}\) is the \defi{hyperbolic reduction} of \(q\colon\Sym^2 \cal E\to\cal L^\vee\) with respect to the regular isotropic subbundle \(\cal F\). We~also say that \(\overline{\cal Q}\coloneqq(\overline{q}=0)\) is the \defi{hyperbolic reduction} of \(\cal Q=(q=0)\) with respect to \(\bb P_S(\cal F)\).
\end{defn}

The process of hyperbolic reduction along a regular isotropic subbundle preserves the degeneration divisor of a quadric fibration:

\begin{lem}[{\cite[Cor.\,1.3.9]{ABB14}}]\label{lem:degeneration-divisor-preserved}
Let $S$ be a smooth scheme over a field of characteristic $\neq 2$, and $\pi\colon \mathcal{Q}\to S$ be a quadric fibration with smooth generic fiber.
Let $\cal F\subset \cal E$ be a regular isotropic subbundle, and $\overline{\pi}\colon \overline{\cal Q}\to S$ be the hyperbolic reduction of~$\mathcal{Q}$ respect to $\bb P_S(\cal F)$.
Then $\pi$ and $\overline{\pi}$ have the same degeneration divisor.
\end{lem}

Note that, more generally, hyperbolic reduction preserves the locus of corank \(\geq i\) fibers for each \(i\) \cite[Lem.\,2.4]{KuznetsovShinder18}.

Hyperbolic reduction can be described geometrically in terms of the linear projection of \(\cal Q\subset\bb P_S(\cal E)\) from the linear subbundle \(\bb P_S(\cal F)\subset\cal Q\subset\bb P_S(\cal E)\) \cite[Prop.\,2.5]{KuznetsovShinder18}.

\subsection{Lemmas of Reid on pencils of quadrics}

Next, over \emph{algebraically closed} fields, we~recall several results proved by Reid \cite{Reid-thesis} that we will use in the proof of Theorem~\ref{thm:symmetric}. In~what follows, \defi{coordinate points} $p_0,\dots, p_N\in \P^N$ mean a choice of coordinates for $\P^N$ such that the $N+1$ points given by intersecting coordinate hyperplanes are exactly $p_0,\dots, p_N$.

\begin{lem}[{\cite[Lem.\,2.2]{Reid-thesis}}]\label{lem:reid_1}
Over an algebraically closed field of characteristic $\neq 2$, let $X$ be a smooth complete intersection of two quadrics in $\P^N$, and let $\ell$ be an $r$-plane on $X$.
Then there exist coordinate points $p_0,\dots, p_N\in \P^N$ such that $\langle p_0,\dots, p_r\rangle =\ell$ and $X$ is defined by two quadrics which correspond in these coordinates to symmetric matrices of the form\vspace*{-3pt}
\[
\begin{pmatrix}
0 & I_{r+1} & 0\\
I_{r+1} & 0 & 0\\
0 & 0 & I_{N-2r-1}
\end{pmatrix}, \quad\arraycolsep4.5pt
\begin{pmatrix}
0 & M & *\\
M & 0 & *\\
* & * & *
\end{pmatrix},
\]
where $M$ is a diagonal \((r+1)\times(r+1)\) matrix with distinct diagonal entries.
\end{lem}

\begin{lem}[{\cite[Lem.\,3.4]{Reid-thesis}}]\label{lem:reid_2}
Over an algebraically closed field of characteristic $\neq 2$, let $\ell, m$ be disjoint $r$-planes in $\P^{2r+1}$
and choose coordinate points $p_0,\dots, p_r, q_0,\dots, q_r\in \P^{2r+1}$ such that $\langle p_0,\dots,p_r\rangle=\ell$ and $\langle q_0,\dots, q_r\rangle =m$.
Furthermore, let $Q_1, Q_2$ be quadrics in $\P^{2r+1}$ which correspond in these coordinates to symmetric matrices of the form
\[
\begin{pmatrix}
0 & I_{r+1}\\
I_{r+1} & 0
\end{pmatrix}, \quad
\begin{pmatrix}
0 & M\\
M & 0
\end{pmatrix},
\]
where $M$ is a diagonal \((r+1)\times(r+1)\) matrix with distinct diagonal entries.
Then the set of $r$-planes on the singular complete intersection of quadrics $Y \coloneqq Q_1\cap Q_2$ coincides with
\[
\left\{\langle p_i, q_j\rangle_{i\in I, j\not\in I}\mid I\subset \left\{0,\dots, r\right\}
\right\}.
\]
\end{lem}

\begin{lem}\label{lem:reid-3}
Over an algebraically closed field of characteristic $\neq 2$,
fix $g\geq 1$, and let $X$ be a smooth complete intersection of two quadrics in $\P^{2g}$.
Let $\ell, m$ be disjoint $(g-1)$-planes on $X$.
Then there exist coordinate points $p_0,\dots, p_{2g}\in \P^{2g}$ such that $\langle p_0,\dots, p_{g-1}\rangle =\ell$, $\langle p_{g},\dots, p_{2g-1}\rangle=m$, and $X$ is defined by two quadrics which correspond in these coordinates to symmetric matrices of the form
\[\arraycolsep4.5pt
\begin{pmatrix}
0 & I_g & *\\
I_g & 0 & *\\
* & * & *
\end{pmatrix}, \quad
\begin{pmatrix}
0 & M & *\\
M & 0 & *\\
* & * & *
\end{pmatrix},
\]
where $M$ is a diagonal \(g\times g\) matrix with distinct diagonal entries.
\end{lem}

For the proof, we~need the following result stated in \cite[p.\,44]{Reid-thesis}:

\begin{lem}\label{lem:elementary}
Over an algebraically closed field of characteristic $\neq 2$, let $A, B$ be $n\times n$ matrices.
Then the following conditions are equivalent:
\begin{enumerate}
\item\label{item:change-of-basis} there exists an invertible $2n\times 2n$ matrix $L$ of the form
\[
L = \begin{pmatrix}
L_1 & 0\\
0 & L_2
\end{pmatrix}
\]
for some $n\times n$ matrices $L_1,L_2$,
and there exists a diagonal $n\times n$ matrix $M$ with distinct diagonal entries such that
\[
L^T
\begin{pmatrix}
0 & A \\
A^T & 0 \\
\end{pmatrix}L =
\begin{pmatrix}
0 & I_n \\
I_n & 0 \\
\end{pmatrix}, \quad
L^T
\begin{pmatrix}
0 & B \\
B^T & 0 \\
\end{pmatrix}L =
\begin{pmatrix}
0 & M \\
M & 0 \\
\end{pmatrix};
\]
\item\label{item:polynomial-separable} the polynomial $\det(\lambda A+ B)$ has $n$ distinct roots.
\end{enumerate}
\end{lem}

\begin{proof}
$\eqref{item:change-of-basis}\Rightarrow\eqref{item:polynomial-separable}$
is immediate.
As for $\eqref{item:polynomial-separable}\Rightarrow\eqref{item:change-of-basis}$,
the assumption implies that $A$ is invertible and $BA^{-1}$ has $n$ distinct eigenvalues.
Let $C$ be an invertible $n\times n$ matrix such that $C^{-1}BA^{-1}C$ equals some diagonal $n\times n$ matrix $M$ with distinct diagonal entries.
Then we may take
\[
L \coloneqq \begin{pmatrix}(C^T)^{-1} & 0\\
0& A^{-1}C\end{pmatrix}.\qedhere
\]
\end{proof}
\begin{proof}[Proof of Lemma \ref{lem:reid-3}]
The proof is outlined in
\cite[p.\,44]{Reid-thesis}.
Here is a detailed argument.
Take coordinate points $p_0,\dots, p_{2g}\in \P^{2g}$ such that $\langle p_0,\dots,p_{g-1}\rangle =\ell$ and $\langle p_g,\dots,p_{2g-1}\rangle =m$. In~these coordinates, any two quadrics defining $X$ correspond to symmetric matrices of the form
\[\arraycolsep4.5pt
S=
\begin{pmatrix}
0 & A & *\\
A^T & 0 & *\\
* & * & *
\end{pmatrix}, \quad
T=\begin{pmatrix}
0 & B & *\\
B^T & 0 & *\\
* & * & *
\end{pmatrix},
\]
where $A, B$ are $g\times g$ matrices,
and, by~the smoothness of $X$, we~may choose those quadrics so that $S$ is invertible.
By \cite[Cor.\,3.7]{Reid-thesis}, the polynomial $\det(\lambda A+B)$ divides the polynomial $\det(\lambda S+ T)$,
where the latter has distinct $2g+1$ roots by the smoothness of $X$ \cite[Prop.\,2.1]{Reid-thesis}.
Hence $\det(\lambda A+ B)$ has distinct $g$ roots, which implies that after a suitable coordinate change
we may take $A=I_g$ and $B$ to be diagonal with distinct diagonal entries by Lemma \ref{lem:elementary}.
Accordingly, we~obtain coordinate points with the desired properties.
\end{proof}

\section{Fano schemes and hyperbolic reductions of pencils of quadrics}
\subsection{Construction of $\phi^{(r)}\colon\mathcal{Q}^{(r)}\to\P^1$ and its properties}\label{sec:Q^(r)-construction}

Throughout Section~\ref{sec:Q^(r)-construction}, we~will work in the following setting.
Fix an arbitrary field $k$ of characteristic $\neq 2$ and an integer $N\geq 2$. Let $X$ be a smooth complete intersection of two quadrics in $\P^N$, and let $\phi\colon \mathcal{Q}\to \P^1$ be the associated pencil of quadrics.
Let $r$ be such that $0\leq r\leq \lfloor\sfrac{N}{2}\rfloor -1$.
By Lemma~\ref{lem:Kuznetsov_fanoscheme}, the Fano scheme $F_r(X)$ of $r$-planes on $X$ is non-empty. We~will assume $F_r(X)(k)\neq \emptyset$ in what follows.

Choose
$\ell\in F_r(X)(k)$. The projection $\pi_{\ell}\colon\bb P^N \dashrightarrow \bb P^{N-r-1}$ away from $\ell$ induces diagrams
\[
\begin{array}{c}
\begin{tikzcd}[column sep=tiny]
&\P_{\P^{N-r-1}}(\O^{\oplus r+1}\oplus \O(1))\arrow[ld, "\ts \bl_{\ell}"'] \arrow[rd] & \\
\P^{N} \arrow[rr, "\ts \pi_{\ell}", dashed]& & \P^{N-r-1},
\end{tikzcd}
\quad
\begin{tikzcd}
& \widetilde{X} \arrow[ld, "\ts \bl_{\ell}"'] \arrow[rd, "\ts f"] & \\
X \arrow[rr, "\ts \pi_{\ell}|_X", dashed]& & \P^{N-r-1},
\end{tikzcd}
\\
\begin{tikzcd}[column sep=scriptsize]
& \widetilde{\mathcal{Q}} \arrow[ld, "\ts \bl_{\P^1\times \ell}"'] \arrow[rd, "\ts h"] & \\
\mathcal{Q} \arrow[rr, "\ts \id\times \pi_{\ell}|_{\mathcal{Q}}", dashed]& & \P^1\times \P^{N-r-1}.
\end{tikzcd}
\end{array}
\]

To be more explicit, choose homogeneous coordinates $x_0,\dots, x_{N}$ for $\P^{N}$ so that
\[
\ell=\bigl\{x_{r+1}=\dots= x_{N}=0\bigr\}.
\]
Then there exist forms $l_{ij}, q_i\in k[x_{r+1}, \dots, x_{N}]$ with $\deg l_{ij}=1, \deg q_i =2$ such that
\begin{align*}
X
=\left\{
\begin{pmatrix}
l_{00} & \dots &l_{0r} & q_0\\
l_{10} & \dots & l_{1r} & q_1
\end{pmatrix}
\begin{pmatrix}
x_0\\[-3pt]
\vdots\\[-3pt]
x_r\\
1
\end{pmatrix}
=0
\right\}
\subset \P^{N}.
\end{align*}
Choosing a section $z$ whose zero set equals the divisor
\[
\P_{\P^{N-r-1}}(\O^{\oplus r+1})\subset \P_{\P^{N-r-1}}(\O^{\oplus r+1 }\oplus \O(1)),
\]
there exist homogeneous coordinates $y_{r+1},\dots, y_{N}$ for $\P^{N-r-1}$ such that $x_l = y_l z$,
and
\[
\widetilde{X}=
\left\{
\begin{pmatrix}
l_{00} & \dots &l_{0r} & q_0\\
l_{10} & \dots & l_{1r} & q_1
\end{pmatrix}
\begin{pmatrix}
x_0\\[-3pt]
\vdots\\[-3pt]
x_r\\
z
\end{pmatrix}
=0
\right\}\subset \P_{\P^{N-r-1}}(\O^{r+1}\oplus \O(1)),
\]
where $l_{ij}, q_i$ are in $y_{r+1},\dots, y_{N}$. We~also have
\begin{align*}
&\mathcal{Q}
=\left\{
\begin{pmatrix}
s & t
\end{pmatrix}
\begin{pmatrix}
l_{00} & \dots &l_{0r} & q_0\\
l_{10} & \dots & l_{1r} & q_1
\end{pmatrix}
\begin{pmatrix}
x_0\\[-3pt]
\vdots\\[-3pt]
x_r\\
1
\end{pmatrix}
=0
\right\}
\subset \P^1\times \P^{N},\\
&\widetilde{\mathcal{Q}}=
\left\{
\begin{pmatrix}
s & t
\end{pmatrix}
\begin{pmatrix}
l_{00} & \dots &l_{0r} & q_0\\
l_{10} & \dots & l_{1r} & q_1
\end{pmatrix}
\begin{pmatrix}
x_0\\[-3pt]
\vdots\\[-3pt]
x_r\\
z
\end{pmatrix}
=0
\right\}\subset \P^1\times \P_{\P^{N-r-1}}(\O^{\oplus r+1}\oplus \O(1)).
\end{align*}

We now define
\begin{align*}
&\mathcal{P}^{(r)}\coloneqq
\left\{
\begin{pmatrix}
s & t
\end{pmatrix}
\begin{pmatrix}
l_{00} & \dots &l_{0r}\\
l_{10} & \dots & l_{1r}
\end{pmatrix}
=0
\right\}
\subset \P^1\times \P^{N-r-1}, \\
&\mathcal{Q}^{(r)}
\coloneqq
\left\{
\begin{pmatrix}
s & t
\end{pmatrix}
\begin{pmatrix}
l_{00} & \dots &l_{0r} & q_0\\
l_{10} & \dots & l_{1r} & q_1
\end{pmatrix}
=0
\right\}
\subset \P^1\times \P^{N-r-1}.
\end{align*}
The first projection defines a $\P^{N-2r-2}$-bundle $\mathcal{P}^{(r)}\to \P^1$, which restricts to a morphism $\phi^{(r)}\colon \mathcal{Q}^{(r)}\to \P^1$, which is a quadric fibration if \(\dim\mathcal Q^{(r)} \geq 1\).
(By convention, $\mathcal{Q}^{(-1)}=\mathcal{Q}$ and $\phi^{(-1)}=\phi$.)
Furthermore, define
\[
E^{(r)}\coloneqq h^{-1}(\mathcal{Q}^{(r)})=\P_{\mathcal{Q}^{(r)}}(\O^{\oplus r+1}\oplus \O(0,1)),
\]
and let $\pi^{(r)}\colon E^{(r)}\to \mathcal{Q}^{(r)}$ be the projection.

\begin{lem}\label{lem:hyperbolic-reduction}
$\mathcal{Q}^{(r)}$ satisfies the following properties.
\begin{enumerate}
\item\label{item:hyperbolic-reduction-pencil} $\phi^{(r)}\colon\mathcal{Q}^{(r)}\to \P^1$
is the hyperbolic reduction of $\phi\colon \mathcal{Q}\to \P^1$ with respect to $\ell$. In~particular, the degeneracy locus of $\phi^{(r)}$ is defined by a separable polynomial of degree $N+1$.
\item\label{item:hyperbolic-reduction-smooth} $\mathcal{Q}^{(r)}$ is smooth of dimension $N-2r-2$. If~$N-2r-2>0$, $\mathcal{Q}^{(r)}$ is geometrically connected.
\item\label{item:hyperbolic-reduction-embedding} $\pi^{(r)}\colon E^{(r)}\to \mathcal{Q}^{(r)}$ induces an embedding $\mathcal{Q}^{(r)}\hto F_{r+1}(\mathcal{Q}/\P^1)$ over $\P^1$,
whose image is the relative Fano
scheme of isotropic $(r+1)$-planes
of
$\phi\colon \mathcal{Q}\to \P^1$
containing~$\ell$.
\item\label{item:hyperbolic-reduction-image}
The restriction
$\pi^{(r)}|_{E^{(r)}\cap (\P^1\times \widetilde{X})}\colon E^{(r)}\cap (\P^1\times \widetilde{X})\to \mathcal{Q}^{(r)}$ induces an isomorphism\vspace*{-5pt}
\begin{multline*}
\mathcal{Q}^{(r)}\setminus\bigl\{l_{ij} = 0 \, (i= 0,1, j=1,\dots, r)\bigr\}\\
\xrightarrow{\,\sim\,} \bigl\{m\in F_r(X)\mid \dim \ell\cap m = r-1 \text{ and }\langle \ell,m\rangle \not\subset X\bigr\}.
\end{multline*}
\item\label{item:hyperbolic-reduction-independence}
If $N\!-\!2r\!-\!2\!>\!0$, the $k$\nobreakdash-birational equivalence class of $\mathcal{Q}^{(r)}$
does not depend on~$\ell$.
\item\label{item:hyperbolic-reduction-max-min} If $N>2$, $\mathcal{Q}^{(0)}$ is $k$\nobreakdash-birational to $X$. If~$N=2g$ for some $g\geq 1$, $\mathcal{Q}^{(g-1)}$ is of dimension zero and length $2g+1$. If~$N=2g+1$ for some $g\geq 1$, $\mathcal{Q}^{(g-1)}$ is $k$\nobreakdash-isomorphic to the curve $C$ of genus $g$ obtained as the Stein factorization of $F_{g}(\mathcal{Q}/\P^1)\to \P^1$.
\item\label{item:index}
If $N-2r-2>0$, $\mathcal{Q}^{(r)}$ has a $0$-cycle of degree $1$,
in other words, the index of $\mathcal{Q}^{(r)}$ is $1$.
\item\label{item:rationalpoint}
If $N-3r-3\geq 0$, $\mathcal{Q}^{(r)}$ has a $k$\nobreakdash-point.
\end{enumerate}
\end{lem}

\skpt
\begin{proof}
\eqref{item:hyperbolic-reduction-pencil}
For $x =[s:t]\in \P^1$, $\mathcal{Q}_{\kappa(x)}$ corresponds to the symmetric matrix
\[
\begin{pmatrix}
0 & A\\
A^T & B
\end{pmatrix},
\]
where $A$ is the Jacobian matrix for $\{s l_{00} + t l_{10} = \dots = s l_{0r} + t l_{1r} = 0\}$ and $B$ is the Hessian matrix for $s q_0 + t q_1$.
This in turn shows that
\[
\ell^{\perp}=
\left\{
\begin{pmatrix}
s & t
\end{pmatrix}
\begin{pmatrix}
l_{00} & \dots &l_{0r}\\
l_{10} & \dots & l_{1r}
\end{pmatrix}
=0
\right\}
\subset \P^{N}_{\kappa(x)}
\]
and the hyperbolic reduction of $\mathcal{Q}_{\kappa(x)}$ with respect to $\ell$ equals $\mathcal{Q}^{(r)}_{\kappa(x)} \subset \mathcal{P}^{(r)}_{\kappa(x)}=\ell^\perp/\ell \subset \P^{N-r-1}_{\kappa(x)}=\P^{N}_{\kappa(x)}/\ell$.
The first statement follows.
As for the second statement, Lemma~\ref{lem:degeneration-divisor-preserved} shows that the degeneracy locus of $\phi^{(r)}\colon \mathcal{Q}^{(r)}\to \P^1$ is the same as that of $\phi\colon \mathcal{Q}\to \P^1$, and the latter is defined by a separable polynomial of degree $N+1$ by \cite[Prop.\,2.1]{Reid-thesis}.

\eqref{item:hyperbolic-reduction-smooth}
By \eqref{item:hyperbolic-reduction-pencil} and Lemma~\ref{lem:quadric-bundle-simple-degeneration},
$\mathcal{Q}^{(r)}$ is smooth.
Moreover, we~deduce from \eqref{item:hyperbolic-reduction-pencil} that we have $\dim \mathcal{Q}^{(r)}=N-2r-2$.
This implies that $\mathcal{Q}^{(r)}$ is a complete intersection of $r+2$ ample divisors in $\P^1\times \P^{N-r-1}$.
So if $N-2r-2>0$, then $\mathcal{Q}^{(r)}$ is geometrically connected.

\eqref{item:hyperbolic-reduction-embedding}--\eqref{item:hyperbolic-reduction-image}
Using the equations for $\mathcal{Q}$ and $\mathcal{Q}^{(r)}$, we~may observe that
$\pi^{(r)}\colon E^{(r)}\to \mathcal{Q}^{(r)}$ is a $\P^{r+1}$-bundle
whose fibers are
the $(r+1)$-planes of $\phi\colon \mathcal{Q}\to \P^1$
containing $\ell$.
The induced morphism $\mathcal{Q}^{(r)}\to F_{r+1}(\mathcal{Q}/\P^1)$ over $\P^1$ is an isomorphism onto its image
because the blow-up of $\mathcal{Q}$ along $\P^1\times \ell$
yields the inverse map.

Similarly, using the equations for $X$ and $\mathcal{Q}^{(r)}$,
it is direct to see that
\[
X\cap \left\{l_{ij}=0\, (i=0,1, j=0,\dots, r)\right\}=\left\{l_{ij}=q_i = 0 \, (i=0,1, j=0,\dots, r)\right\}\subset \P^N
\]
is the union of all linear spaces on $X$ containing $\ell$,
and that
over the complement of the proper closed subscheme $\mathcal{Q}^{(r)}\cap\left\{l_{ij} = 0 \, (i= 0,1, j=1,\dots, r)\right\}\subset \mathcal{Q}^{(r)}$, the morphism $\pi^{(r)}|_{E^{(r)}\cap (\P^1\times \widetilde{X})}\colon E^{(r)}\cap (\P^1\times \widetilde{X})\to \mathcal{Q}^{(r)}$ is
a $\P^{r}$-bundle whose fibers are the $r$-planes $m$ on $X$ such that $\dim \ell\cap m = r-1$ and $\langle \ell, m\rangle \not\subset X$.
The induced morphism
\begin{multline*}
\mathcal{Q}^{(r)}\setminus\bigl\{l_{ij} = 0 \, (i= 0,1, j=1,\dots, r)\bigr\}\\
\to\bigl\{m\in F_r(X)\mid \dim \ell\cap m = r-1 \text{ and }\langle \ell,m\rangle \not\subset X\bigr\}
\end{multline*}
is an isomorphism because the blow-up of $X$ along $\ell$ yields the inverse map.

\eqref{item:hyperbolic-reduction-independence}
It is enough for us to show that the $k(\P^1)$-isomorphism class of the generic fiber of $\mathcal{Q}^{(r)}\to \P^1$ does not depend on $\ell$.
For $\ell, \ell'\in F_r(X)(k)$, denote by $\mathcal{Q}^{(r)}_\ell\to \P^1$, $\mathcal{Q}^{(r)}_{\ell'}\to \P^1$
the corresponding hyperbolic reductions. Letting $H$ denote a split quadric, we~have
\[
Q_{k(\P^1)}\simeq (\mathcal{Q}^{(r)}_{\ell})_{k(\P^1)}\perp H \simeq (\mathcal{Q}^{(r)}_{\ell'})_{k(\P^1)}\perp H.
\]
The Witt cancellation theorem \cite[Th.\,8.4]{EKM-quadratic-forms}
then shows $(\mathcal{Q}^{(r)}_\ell)_{k(\P^1)}\simeq (\mathcal{Q}^{(r)}_{\ell'})_{k(\P^1)}$.

\eqref{item:hyperbolic-reduction-max-min}
If $N>2$, then $\mathcal{Q}^{(0)}$ is $k$\nobreakdash-birational to the image of $X$ under the projection away from a $k$\nobreakdash-point on $X$, hence $k$\nobreakdash-birational to $X$. If~$N=2g$ for some $g\geq 1$, $\mathcal{Q}^{(g-1)}$ is a complete intersection of $g$ divisors of type $(1,1)$ and a divisor of type $(1,2)$ in $\P^1\times \P^{g}$, hence $\mathcal{Q}^{(g-1)}$ is of dimension zero and length $2g+1$. If~$N=2g+1$ for some $g\geq 1$,
the composition $\mathcal{Q}^{(g-1)}\hto F_{g}(\mathcal{Q}/\P^1)\to C$ is an isomorphism over \(\P^1\)
since the double covers $\phi^{(g-1)}\colon\mathcal{Q}^{(g-1)}\to \P^1$ and $C\to \P^1$ are both branched over the degeneracy locus of $\phi\colon \mathcal{Q}\to \P^1$.

\eqref{item:index}
Since $\mathcal{Q}^{(r)}$ is a complete intersection of $r+1$ divisors of type $(1,1)$ and a divisor of type $(1,2)$ in $\P^1\times \P^{N-r-1}$,
the zero-cycle $\mathcal{Q}^{(r)}\cdot (\P^1\times \P^{r+1}) - (r+1) \mathcal{Q}^{(r)}\cdot (*\times \P^{r+2})$ is of degree $1$.

\eqref{item:rationalpoint}
If $N-3r-3\geq 0$,
then $\left\{l_{ij}=0\, (i=0,1, j=1,\dots, r)\right\} \subset\P^{N-r-1}$
is non-empty by a dimension count.
Using the equations for $\mathcal{Q}^{(r)}$
and the fact that $l_{ij}$ are linear forms, $\mathcal{Q}^{(r)}\cap\left\{l_{ij}=0\, (i=0,1, j=0,\dots, r)\right\}$ has a $k$\nobreakdash-point.
This concludes the proof.
\end{proof}

\subsection{Complete original proof of Theorem~\ref{thm:symmetric}}
\begin{proof}[Proof of Theorem~\ref{thm:symmetric}]
Under the assumptions of Theorem~\ref{thm:symmetric}, choose $\ell\in F_r(X)(k)$ so that $\mathcal{Q}^{(r)}$ is defined.
For $m\in F_r(X)$, not necessarily a $k$\nobreakdash-point, we~say that $m$ is \defi{general with respect to $\ell$} if:
(1) $\ell\cap m=\emptyset$; and (2) $X\cap \langle \ell, m\rangle \subset \langle \ell, m\rangle = \P^{2r+1}$ is defined by two quadrics which, over an algebraic closure of the residue field of $m$ and for some choice of coordinates, correspond to symmetric matrices of the form
\[
\begin{pmatrix}
0 & I_{r+1}\\
I_{r+1} & 0
\end{pmatrix}, \quad
\begin{pmatrix}
0 & M\\
M & 0
\end{pmatrix},
\]
where $M$ is diagonal with distinct diagonal entries.
By Lemma \ref{lem:elementary},
property (2) is equivalent to separability of a certain polynomial associated to $\ell$ and $m$. Define
\[
U\coloneqq\bigl\{m \in F_r(X) \mid m \text{ is general with respect to }\ell \bigl\}.
\]
The set $U$ is a priori only defined over $\kbar$, but it actually descends to $k$ as $\ell$ is defined over $k$.
Clearly, $U$ is open in $F_r(X)$, and moreover, it is non-empty by Lemma~\ref{lem:reid_1}.
For every $m\in U$, Lemma~\ref{lem:reid_2} implies that the (singular) intersection $X\cap \langle \ell, m\rangle \subset \langle \ell,m\rangle = \bb P^{2r+1}$ contains exactly $r+1$ distinct $r$-planes $m_1,\dots, m_{r+1}$ that intersect~$\ell$ along an $(r-1)$-plane.
Using the isomorphism in Lemma \ref{lem:hyperbolic-reduction}\eqref{item:hyperbolic-reduction-image}, we~now define a morphism
\[
\psi\colon U\to \Sym^{r+1}\mathcal{Q}^{(r)},\quad
m\mto (m_1,\dots, m_{r+1}).
\]
The morphism $\psi$ is defined over $k$ by the same reasoning as above.
Moreover, $\psi$ is one-to-one onto its image on the level of $\kbar$-points
because we have $\langle m_1,\dots, m_{r+1}\rangle =\langle \ell,m\rangle$ and $m$ is the unique $r$-plane on $X\cap \langle \ell,m\rangle$ that is disjoint from $\ell$.
Since $\dim F_r(X) = \dim \Sym^{r+1}\mathcal{Q}^{(r)}=(r+1)(N-2r-2)$ by Lemma \ref{lem:hyperbolic-reduction}\eqref{item:hyperbolic-reduction-smooth}
and $\psi$ factors through the smooth locus of $\Sym^{r+1}\mathcal{Q}^{(r)}$,
it then follows that
$\psi$ is an open immersion. If~\hbox{$N-2r-2>0$}, then $F_r(X)$ and $\Sym^{r+1}\mathcal{Q}^{(r)}$ are both geometrically integral by Lemmas~\ref{lem:Kuznetsov_fanoscheme} and~\ref{lem:hyperbolic-reduction}, so $F_r(X)$ is $k$\nobreakdash-birational to $\Sym^{r+1}\mathcal{Q}^{(r)}$.

It remains to consider the case $N-2r-2=0$. Set $g=r+1$. We~have $\dim F_{g-1}(X)=\dim \Sym^{g}\mathcal{Q}^{(g-1)}=0$. In~this case, Lemma~\ref{lem:reid-3} shows that if $m\in F_{g-1}(X)$ satisfies property (1) of being general with respect to \(\ell\), then property (2) is automatic.
Hence~$U$ coincides with the subscheme of $F_{g-1}(X)$ parametrizing $(g-1)$-planes on $X$ disjoint from $\ell$.
It remains to show that $\psi$ gives an isomorphism between $U$ and the subscheme of $\Sym^{g}\mathcal{Q}^{(g-1)}$ parametrizing $g$-tuples of distinct points of $\mathcal{Q}^{(g-1)}$.
The former is of length $\binom{2g+1}{g}$ by the analysis on the configuration of the $(g-1)$-planes on $X$ due to Reid \cite[Th.\,3.8]{Reid-thesis},
and the latter is also of length $\binom{2g+1}{g}$
by Lemma~\ref{lem:hyperbolic-reduction}\eqref{item:hyperbolic-reduction-max-min}. The claim thus follows, and this concludes the proof of the theorem.
\end{proof}

\begin{thebibliography}{EKM08}
\bibitem[ABB14]{ABB14} A. Auel, M. Bernardara and M. Bolognesi, ``Fibrations in complete intersections of quadrics, Clifford algebras, derived categories, and rationality problems'', \emph{J. Math. Pures Appl.} (9) \textbf{102} (2014), no. 1, 249--291.
\bibitem[CT24]{CT-intersection-quadrics} J.-L. Colliot-Th\'el\`ene, ``Retour sur l'arithm\'etique des intersections de deux quadriques'', \emph{J. reine angew. Math.} \textbf{806} (2024), 147--185, avec un appendice par A. Kuznetsov.
\bibitem[EKM08]{EKM-quadratic-forms} R. Elman, N. Karpenko and A. Merkurjev, \emph{The algebraic and geometric theory of quadratic forms}, Amer. Math. Soc. Colloquium Publ. 56, American Mathematical Society, Providence, RI, 2008.
\bibitem[Kuz24]{Kuznetsov-hyperbolic} A. Kuznetsov, ``Quadric bundles and hyperbolic equivalence'', \emph{Geom. Topol.} \textbf{28} (2024), no. 3, 1287--1339.
\bibitem[KS18]{KuznetsovShinder18} A. Kuznetsov and E. Shinder, ``Grothendieck ring of varieties, D- and L-equivalence, and families of quadrics'', \emph{Selecta Math. (N.S.)} \textbf{24} (2018), no. 4, 3475--3500.
\bibitem[Rei72]{Reid-thesis} M. Reid, ``The complete intersection of two or more quadrics'', PhD Thesis, Trinity College, Cambridge, 1972, \url{https://mreid.warwick.ac.uk/3folds/qu.pdf}.
\end{thebibliography}
\end{document}
